\subsection{Follow-up interventions}
\label{sec:followups}

The analyses above suggest three explanations, each with a direct test. We train each intervention with the monitorability-aware reward, invariant observations and 10 seeds, and compare it with the reference agent on paired episodes (Table~\ref{tab:followups}). These comparisons are exploratory and their $p$-values are not corrected for multiplicity.

\begin{table}[t]
\centering
\small
\setlength{\tabcolsep}{3pt}
\caption{Follow-up interventions (mean over 10 seeds, 95\% $t$-interval). ``Last ckpt.'' evaluates the last checkpoint instead of the selected one; ``collapsed'' is the fraction of late checkpoints (after 5{,}000 episodes) with training-length success $\le 0.35$; ``wasted \texttt{REVISE}'' counts non-repair \texttt{REVISE}s per episode at lengths 10--14.}
\label{tab:followups}
\resizebox{\linewidth}{!}{\begin{tabular}{lcccccc}
\toprule
Variant & Success 3--6 & Success 7--9 & Success 10--14 & Last ckpt.\ 7--9 & Collapsed & Wasted REVISE \\
\midrule
Monitorability (inv) & 0.91\,{\scriptsize[0.89, 0.92]} & 0.86\,{\scriptsize[0.83, 0.89]} & 0.74\,{\scriptsize[0.68, 0.80]} & 0.63\,{\scriptsize[0.51, 0.76]} & 21\% & 0.74 \\
+ slack-per-step feature & 0.83\,{\scriptsize[0.80, 0.87]} & 0.46\,{\scriptsize[0.35, 0.57]} & 0.28\,{\scriptsize[0.20, 0.35]} & 0.37\,{\scriptsize[0.31, 0.43]} & 38\% & 0.64 \\
+ fault-located memory bit & 0.94\,{\scriptsize[0.92, 0.95]} & 0.88\,{\scriptsize[0.86, 0.91]} & 0.74\,{\scriptsize[0.65, 0.83]} & 0.81\,{\scriptsize[0.70, 0.91]} & 10\% & 0.23 \\
fault bonus paid on repair & 0.89\,{\scriptsize[0.88, 0.91]} & 0.79\,{\scriptsize[0.71, 0.87]} & 0.64\,{\scriptsize[0.52, 0.75]} & 0.75\,{\scriptsize[0.67, 0.83]} & 1\% & 0.51 \\
memory bit + bonus on repair & 0.93\,{\scriptsize[0.92, 0.95]} & 0.83\,{\scriptsize[0.76, 0.91]} & 0.66\,{\scriptsize[0.54, 0.77]} & 0.79\,{\scriptsize[0.72, 0.85]} & 0\% & 0.15 \\
+ repair-pending bit & 0.95\,{\scriptsize[0.94, 0.95]} & 0.83\,{\scriptsize[0.76, 0.91]} & 0.60\,{\scriptsize[0.45, 0.75]} & 0.89\,{\scriptsize[0.87, 0.91]} & 0\% & 0.08 \\
\bottomrule
\end{tabular}
}
\end{table}

\paragraph{Remaining-length information in unseen combinations does not help.}
Adding slack per remaining step, $\mathrm{clip}(\mathrm{slack}/(d-t), 0, 2)/2$, gives the agent the information that \S\ref{sec:residual} identifies as missing. Its values stay inside the interval seen in training, but many of them, and most of the observation vectors they form, never occur on training lengths. Observations at \val{coverage/ratio/heldout_long/unseen_rate/pct} of decision points at lengths 7--9 and \val{coverage/ratio/heldout_xlong/unseen_rate/pct} at 10--14 never appear in the training-length reference sample, against \val{coverage/invariant/heldout_long/unseen_rate/pct} and \val{coverage/invariant/heldout_xlong/unseen_rate/pct} for the invariant encoder. Held-out success falls to \val{followups/ma_ratio/heldout_long/success} at 7--9 and \val{followups/ma_ratio/heldout_xlong/success} at 10--14 ($\val{diff/ma_ratio-ma_inv/heldout_xlong/success}$, $p=\val{diff/ma_ratio-ma_inv/heldout_xlong/success/p}$). Fault detection stays high (\val{followups/ma_ratio/heldout_xlong/fault_detected}), but repairs are rarely completed (\val{followups/ma_ratio/heldout_xlong/success_given_detected}). The extra feature also hurts in distribution, by $\val{diff/ma_ratio-ma_inv/train_like/success}$ on training lengths, so its effect cannot be attributed to out-of-range inputs alone. The result is consistent with the coverage explanation, and it shows that adding the missing information is not enough.

\paragraph{A memory bit removes most wasted \texttt{REVISE}s.}
Adding one bit that records whether an \textsc{unsafe} check has been seen removes the aliasing between repairs and no-ops. Wasted \texttt{REVISE}s at 10--14 fall from \val{main/ma_inv/heldout_xlong/nonrepair_revise_count} to \val{followups/ma_memory/heldout_xlong/nonrepair_revise_count} per episode. Training-length success rises by $+$\val{diff/ma_memory-ma_inv/train_like/success} (\val{diff/ma_memory-ma_inv/train_like/success/npos}/10 seeds, $p=\val{diff/ma_memory-ma_inv/train_like/success/p}$). Held-out success at selected checkpoints does not change resolvably ($p=\val{diff/ma_memory-ma_inv/heldout_xlong/success/p}$ at 10--14), because the agent spends the recovered slack on more un-cued checks (\val{followups/ma_memory/heldout_xlong/verify_rate_cue0} vs.\ \val{main/ma_inv/heldout_xlong/verify_rate_cue0}). The allocation explanation of \S\ref{sec:residual} is therefore only partly confirmed: removing one source of waste was not enough. Checkpoint collapse is lower with the memory bit (\val{followups/ma_memory/collapse_rate/pct} vs.\ \val{main/ma_inv/collapse_rate/pct}), but the intervals overlap. The bit changes both sides of the probe. Continuing past a caught fault becomes less costly in truth (true gap \val{probe/ma_memory/true_gap}, against \val{probe/ma_inv/true_gap} without the bit), because these agents often go back and repair the fault later (a return of \val{probe/ma_memory/return_think} after continuing, against \val{probe/ma_inv/return_think}; success after continuing requires undoing back to the fault). The learned gap stays small (\val{probe/ma_memory/q_gap}).

\paragraph{Paying at the repair step removes the detect-without-repair collapse.}
Moving the $0.25$ fault bonus from detection to the first correct re-execution of the located fault places a reward at the repair step itself. If the collapse of \S\ref{sec:norepair} is a credit-assignment failure, this should help. The test changes two things at once: it removes the $0.25$ detection payment and adds a repair payment, and we do not separate them. The cue bonus is still paid when a cued fault is caught. The ablation that removes the detection payment (together with the safe-check bonus) \emph{without} adding a repair payment (no check bonuses, \S\ref{sec:ablations}) still collapses in \val{ablations/no_check_bonuses/collapse_rate/pct} [\val{ablations/no_check_bonuses/collapse_rate/lo/pct}, \val{ablations/no_check_bonuses/collapse_rate/hi/pct}] of late checkpoints. The repair payment, not the missing detection payment, is therefore the likely active ingredient. Collapsed late checkpoints drop from \val{main/ma_inv/collapse_rate/pct} [\val{main/ma_inv/collapse_rate/lo/pct}, \val{main/ma_inv/collapse_rate/hi/pct}] to \val{followups/ma_repair/collapse_rate/pct} [\val{followups/ma_repair/collapse_rate/lo/pct}, \val{followups/ma_repair/collapse_rate/hi/pct}]. The mean absolute change between consecutive evaluations drops from \val{main/ma_inv/mean_abs_change} to \val{followups/ma_repair/mean_abs_change}. With the memory bit added as well, collapse disappears entirely (\val{followups/ma_memory_repair/collapse_rate/pct}). The repair payment widens the learned gap between repairing and continuing after a caught fault, from \val{probe/ma_inv/q_gap} to \val{probe/ma_repair/q_gap}, and to \val{probe/ma_memory_repair/q_gap} with the memory bit (medians over \textsc{unsafe} decision points). The true gaps are about two (\val{probe/ma_inv/true_gap}, \val{probe/ma_repair/true_gap} and \val{probe/ma_memory_repair/true_gap}), so even the larger learned gaps are far too small, but they appear large enough to keep the greedy choice stable.

The benefit shows up in \emph{reliability} rather than peak performance. At selected checkpoints the repair-paid reward does not beat the original, because selection already discards collapsed checkpoints. It may even be slightly worse there: \val{diff/ma_repair-ma_inv/heldout_long/success} [\val{diff/ma_repair-ma_inv/heldout_long/success/lo}, \val{diff/ma_repair-ma_inv/heldout_long/success/hi}] at 7--9 ($p=\val{diff/ma_repair-ma_inv/heldout_long/success/p}$) and \val{diff/ma_repair-ma_inv/heldout_xlong/success} [\val{diff/ma_repair-ma_inv/heldout_xlong/success/lo}, \val{diff/ma_repair-ma_inv/heldout_xlong/success/hi}] at 10--14, differences that 10 seeds do not resolve. At the last checkpoint it is better: $+$\val{diff/ma_repair@final-ma_inv@final/heldout_long/success} [\val{diff/ma_repair@final-ma_inv@final/heldout_long/success/lo}, \val{diff/ma_repair@final-ma_inv@final/heldout_long/success/hi}] at 7--9 ($p=\val{diff/ma_repair@final-ma_inv@final/heldout_long/success/p}$) and $+$\val{diff/ma_repair@final-ma_inv@final/heldout_xlong/success} at 10--14 ($p=\val{diff/ma_repair@final-ma_inv@final/heldout_xlong/success/p}$).

The two last-checkpoint agents behave differently. The repair-paid agent checks less (\val{followups/ma_repair@final/heldout_long/verify_count} vs.\ \val{main/ma_inv@final/heldout_long/verify_count} checks per episode) and catches fewer faults (\val{followups/ma_repair@final/heldout_long/fault_detected} vs.\ \val{main/ma_inv@final/heldout_long/fault_detected}). But it almost always completes the repairs it starts (\val{followups/ma_repair@final/heldout_long/success_given_detected} vs.\ \val{main/ma_inv@final/heldout_long/success_given_detected}).
